\section{Computational details}\label{sec:compdetails}

\statspai{} installs from PyPI without compiled extensions of its own;
its numerical core runs on the scientific \proglang{Python} stack
(\pkg{numba}'s JIT compiler is the only compilation step, and it is
optional at runtime). The submitted artifact uses \proglang{Python}
3.12.1, \pkg{NumPy} 1.26.4,
\pkg{SciPy} 1.13.0, \pkg{pandas} 2.2.0, \pkg{statsmodels} 0.14.1,
\pkg{scikit-learn} 1.4.0, \pkg{linearmodels} 6.0, and
\pkg{matplotlib} 3.8 \citep{hunter2007matplotlib}. Optional heavy backends are isolated under extras
and are lazily imported. The submitted numerical and timing claims do
not require a compiled Rust extension, GPU, or commercial software.

The \proglang{R} parity harness is pinned by
\path{tests/r_parity/renv.lock}. The reference environment uses
\proglang{R} 4.5.2; the packages behind the twelve-estimator suite are
\pkg{fixest} 0.14.0, \pkg{did} 2.3.0,
\pkg{Synth} 1.1-10, \pkg{synthdid} 0.0.9, \pkg{rdrobust} 3.0.0,
\pkg{rddensity} 2.6, \pkg{HonestDiD} 0.2.8, \pkg{DoubleML} 1.0.2,
\pkg{grf} 2.6.1, and \pkg{MatchIt} 4.7.2, joined by \pkg{fect} 2.4.1
and \pkg{interflex} 1.4.0 for the three-ecosystem worked example; the
complete dependency closure for all \RParityModuleCount{} modules
(including \pkg{augsynth}, \pkg{gsynth}, \pkg{lme4}, \pkg{CBPS}
\citep{imai2014covariate}, \pkg{ebal} \citep{hainmueller2012entropy},
\pkg{optmatch} \citep{hansen2006optimal}, \pkg{Matching}
\citep{sekhon2011multivariate}, and \pkg{sbw}
\citep{wang2020minimal}) is recorded
with exact versions and GitHub SHAs in the lockfile and per-row JSON
provenance. The Stata bridge uses \proglang{Stata} 18 MP and records
official/community command versions in
\path{tests/stata_parity/STATA_ENVIRONMENT.md}. All \StataParityModuleCount{} R-joined Stata
modules and all \RParityModuleCount{} \proglang{R} modules reproduce their frozen JSON
outputs with zero drift; the Stata leg requires a licence, so the
archive ships the frozen JSONs, matching do-files, and engine
provenance, and the main paper remains auditable from the
\proglang{Python}/\proglang{R} artifacts.

\subsection{Distribution and installation}

The package is distributed through PyPI
(\code{pip install statspai}); this article describes release \StatsPAIVersion{}
(tag \code{v\StatsPAIVersion{}}). Optional extras are exposed as
\code{dev}, \code{performance}, \code{bayes}, \code{deepiv},
\code{neural}, \code{fixest}, \code{rd-cct}, and \code{plotting}
bundles. The \code{rd-cct} extra installs the official \pkg{rdrobust}
\proglang{Python} port (\code{rdrobust>=2.0}, maintained by the
\pkg{rdrobust} authors), which backs the optional \code{bwselect="cct"}
path of Section~\ref{sec:ex-rd}; that port is GPL-3 licensed, which is
why it is an opt-in extra installed by the user rather than a bundled
dependency of the MIT-licensed core. No parity claim in this paper
depends on it: the native default selector reproduces its numbers on
the audited fixture without it.

The test suite runs on Linux, macOS, and Windows in continuous
integration; the acceptance evidence for this paper is the Track~A--D
artifacts, worked examples, and audit scripts, each regenerated rather
than quoted.

\subsection{Reproducibility of this paper}\label{sec:reproducibility}

Reviewers can rebuild the paper's evidence along three paths, driven by
one standalone script, \path{Paper-JSS/replication/reproduce.py}. The
\emph{licence-free path} (\code{--tier 1}; the default) needs only the
pinned \proglang{Python} environment and runs in under ten minutes on
commodity hardware, inside the JSS one-hour threshold for a reviewer
verification path. It is important to be precise about what it does.
It \emph{recomputes} the seven worked examples, every manuscript listing
and transcript, the known-truth recovery tests, the seed-replicated
forest analysis, and the configuration map; it only \emph{re-derives
tables} from frozen artifacts for the three experiments that cost hours
--- Track~A parity, Track~B Monte Carlo, and Track~C timing. Each of
those has a separate, documented recomputation entry point
(Table~\ref{tab:reproduction}). The \emph{\proglang{R} path} (\code{--tier 2}) adds the
pinned \code{renv} environment and re-executes the cross-language
parity modules themselves, which the licence-free path can only replay
from frozen artifacts; by construction the paper's parity claim is a
comparison against \proglang{R}, so auditing the rows at full strength
requires this path. The \emph{\proglang{Stata} path} (\code{--tier 3})
re-executes the migration bridges and requires a Stata licence. A
Dockerfile reproduces the licence-free environment exactly.

Two further checks are on the paper rather than the package. The
listing audit replays each printed transcript statement by statement and
compares the printed output line for line with what the code produces.
And both sides of every parity row are re-derivable:
\path{tests/r_parity/verify_reproduce.py} re-runs each \proglang{R}
reference on the committed bytes, and
\path{tests/r_parity/verify_reproduce_py.py} does the same for the
\statspai{} side and also regenerates each fixture CSV byte for byte;
all \RParityModuleCount{} modules clear a \(10^{-9}\) reproduction
floor with no exemptions. On the licence-free path the parity,
coverage, and timing experiments are frozen evidence on that path;
Table~\ref{tab:reproduction} lists how to recompute each.

\begin{table}[!ht]
\centering\scriptsize
\begin{tabular}{@{}L{0.19\linewidth}L{0.15\linewidth}L{0.30\linewidth}L{0.13\linewidth}L{0.11\linewidth}@{}}
\toprule
Artifact & Tier~1 does & Full recomputation & Needs & Time \\
\midrule
Worked examples, listings (Sec.~\ref{sec:examples}) & recomputes & \code{reproduce.py} (Tier~1) & \proglang{Python} & \(<\)10 min \\
Track~A, \statspai{} side (Tab.~\ref{tab:track-a-parity}) & re-reads JSON & \code{--recompute parity-py} (and fixture bytes) & \proglang{Python} & \(\sim\)15 min \\
Track~A, \proglang{R} side & re-reads JSON & \code{--tier 2} & \proglang{R} 4.5.2, \code{renv.lock} & 30--60 min \\
Track~A, \proglang{Stata} side & re-reads JSON & \code{--tier 3} & \proglang{Stata} 18 & \(\sim\)30 min \\
Forest seeds (Tab.~\ref{tab:forest-seed-mc}) & re-analyses frozen draws & \code{--recompute forest-seed} & \proglang{Python}; \proglang{R} for \pkg{grf} side & \(\sim\)30 min + \(\sim\)20 min \\
Track~B (Tabs.~\ref{tab:track-b-b1000}--\ref{tab:track-b-robustness}) & re-reads JSON & \code{--recompute montecarlo} & \proglang{Python} & \(\sim\)80 min \\
Track~C (Fig.~\ref{fig:track-c-loglog}) & re-reads JSON & \code{--recompute performance} & \proglang{Python}; \proglang{R} for reference legs & hardware-dependent \\
\bottomrule
\end{tabular}
\caption{What the licence-free path recomputes and how to recompute the
rest (options of \code{replication/reproduce.py}); each writes to the
artifact path the paper reads.}
\label{tab:reproduction}
\end{table} The replication archive contains the
manuscript sources, replication scripts, selected package
source, validation tests, R/Stata/reference parity artifacts, Monte
Carlo and performance artifacts, and the deterministic MCP trace, and
fits within the JSS attachment limit.
